\section*{3 Saving Space}
In these sections, our register programs will all operate over the field of two elements: $R_n=\mathbb F_2$ for all $n$.
\subsection*{3.1 Basic Algorithm}
In this section, we will prove:
\paragraph{Theorem9.} For any function $f:\{0,1\}^n\to\{0,1\}$ there is an $m$-catalytic branching program with length $4n$ and width $2m$ that computes $f$, where $m=2^{n+2^n}$.
This is proved by Algorithm1 below. This nearly reproduces Potechin's Theorem3.1[11], but with a worse value $m=2^{n+2^n}$ instead of $2^{2^n-1}$. Nonetheless, we will find it a useful starting point for our algorithm that trades space for time in Section3.2.
\begin{proof}
We will design a program that uses $n+2^n$ work registers plus one output register $R_{\mathrm{out}}$, which is sufficient by Observation7. First, we have $n$ registers $R^{\mathrm{in}}_1,\ldots,R^{\mathrm{in}}_n$, corresponding to the $n$ input bits. This correspondence is given by the following subroutine:
\paragraph{ToggleInput.}
\begin{enumerate}[label=\arabic*:]
\item \textbf{procedure} ToggleInput
\item \textbf{for} $i=1,\ldots,n$ \textbf{do}
\item $R_i^{\mathrm{in}}\leftarrow R_i^{\mathrm{in}}+x_i$
\item \textbf{end for}
\item \textbf{end procedure}
\end{enumerate}
After ToggleInput runs, the registers have values $R_i^{\mathrm{in}}=\tau_i^{\mathrm{in}}+x_i$, where $\tau_i^{\mathrm{in}}$ stands for the initial value of $R_i^{\mathrm{in}}$. If we run it a second time, the registers are restored to their original values: $R_i^{\mathrm{in}}=\tau_i^{\mathrm{in}}$. Since we query all $n$ variables once, ToggleInput requires time $n$ to run once.
Before defining our other $2^n$ registers, we introduce an algebraic view of $f$, which will be our main focus. We can view $f$ as a multilinear polynomial $p_f\in\mathbb F_2[x_1,\ldots,x_n]$ using basic interpolation:
\paragraph{Lemma10.} For any function $f:\{0,1\}^n\to\{0,1\}$ there is a multilinear polynomial $p_f\in\mathbb F_2[x_1,\ldots,x_n]$ such that $p_f(\vec x)=f(\vec x)$ for all $\vec x\in\{0,1\}^n$.
\begin{proof}
For any $\vec y\in\{0,1\}^n$, we can algebraically define the indicator function $[\vec x=\vec y]$ as $\prod_{i=1}^n(x_i+y_i+1)\in\mathbb F_2[x_1,\ldots,x_n]$. Then we can set
\[p_f=\sum_{\vec y\in\{0,1\}^n:f(\vec y)=1}[\vec x=\vec y].\]
\end{proof}
Now define $y_i=\tau_i^{\mathrm{in}}+x_i$; in other words, $y_i$ is the value of $R_i^{\mathrm{in}}$ after running ToggleInput. Define $q_f(y_1,\ldots,y_n,\tau_1^{\mathrm{in}},\ldots,\tau_n^{\mathrm{in}})=p_f(y_1-\tau_1^{\mathrm{in}},\ldots,y_n-\tau_n^{\mathrm{in}})$ to be the result of rewriting $p_f$ using the $y_i$ and $\tau_i$ variables.\footnote{For example, if $f$ is the OR function $f(x_1,x_2)=x_1\lor x_2$, we have $p_f(x_1,x_2)=x_1+x_2+x_1x_2$ and $q_f(\tau_1^{\mathrm{in}},\tau_2^{\mathrm{in}},y_1,y_2)=y_1+\tau_1^{\mathrm{in}}+y_2+\tau_2^{\mathrm{in}}+y_1y_2+y_1\tau_2^{\mathrm{in}}+\tau_1^{\mathrm{in}}y_2+\tau_1^{\mathrm{in}}\tau_2^{\mathrm{in}}$, and both are equal to $f(x_1,x_2)$ so long as $\vec y$ have the correct values.} For $S,S'\subseteq[n]$, let $c_{S,S'}$ be the coefficient of $(\prod_{i\in S}\tau_i^{\mathrm{in}})\cdot(\prod_{i\in S'}y_i)$ in $q_f$, so that
\[q_f(\vec\tau^{\mathrm{in}},\vec y)=\sum_{S,S'\subseteq[n]}c_{S,S'}\left(\prod_{i\in S}\tau_i^{\mathrm{in}}\right)\left(\prod_{i\in S'}y_i\right).\]
Note that $S$ and $S'$ are disjoint whenever $c_{S,S'}\neq0$, because $p_f$ is multilinear. We now introduce our other registers: we have $2^n$ registers $R_S$ indexed by subsets $S\subseteq[n]$. Our next subroutine prepares us to use these registers to compute $q_f$:
\paragraph{ToggleMonomials.}
\begin{enumerate}[label=\arabic*:]
\item \textbf{procedure} ToggleMonomials
\item ToggleInput
\item \textbf{for} $S'\subseteq[n]$ \textbf{do}
\item $R_{S'}\leftarrow R_{S'}+\prod_{i\in S'}R_i^{\mathrm{in}}$
\item \textbf{end for}
\item ToggleInput
\item \textbf{end procedure}
\end{enumerate}
After ToggleMonomials runs, we have $R_S=\tau_S+\prod_{i\in S}y_i$ for each $S\subseteq[n]$, where $\tau_S$ stands for the register's initial value. The $R^{\mathrm{in}}$ registers have their initial values $R_i^{\mathrm{in}}=\tau_i^{\mathrm{in}}$. We run ToggleInput twice and have $2^n$ additional instructions, but since the additional instructions do not query any $x$ variables they can be computed in the last $x$ query of ToggleInput, for a total runtime of $2n$.
Our final algorithm for computing $f$ is:
\paragraph{Algorithm1: Transparently compute $f$ in time $4n$ with $n+2^n+1$ registers.}
\begin{enumerate}[label=\arabic*:]
\item ToggleMonomials
\item $R_{\mathrm{out}}\leftarrow R_{\mathrm{out}}+\sum_{S,S'\subseteq[n]}c_{S,S'}(\prod_{i\in S}R_i^{\mathrm{in}})R_{S'}$
\item ToggleMonomials
\item $R_{\mathrm{out}}\leftarrow R_{\mathrm{out}}-\sum_{S,S'\subseteq[n]}c_{S,S'}(\prod_{i\in S}R_i^{\mathrm{in}})R_{S'}$
\end{enumerate}
When Line2 executes, we have $R_{S'}=\tau_{S'}+\prod_{i\in S'}y_i$, so after that line,
\[R_{\mathrm{out}}=\tau^{\mathrm{out}}+\sum_{S,S'\subseteq[n]}c_{S,S'}\left(\prod_{i\in S}\tau_i^{\mathrm{in}}\right)\left(\tau_{S'}+\prod_{i\in S'}y_i\right).\]
Then when Line4 executes, we have $R_{S'}=\tau_{S'}$, so the final value is
\begin{multline*}
R_{\mathrm{out}}=\tau^{\mathrm{out}}+\sum_{S,S'\subseteq[n]}c_{S,S'}\left(\prod_{i\in S}\tau_i^{\mathrm{in}}\right)\left(\prod_{i\in S'}y_i\right)\\
=\tau^{\mathrm{out}}+f(x_1,\ldots,x_n).
\end{multline*}
So Algorithm1 correctly computes $f$. The space of the program is $n+2^n$ by construction, and as before we can ignore the instructions on lines2 and4 since they do not use $\mathbf x$, giving us a total runtime of $4n$ from the two calls to ToggleMonomials.
\end{proof}
